\documentclass[a4paper]{article}

% Basic packages
% Español (México) con XeLaTeX: fontspec para fuentes Unicode, polyglossia
% para la división silábica y los nombres del documento en la variante mexicana.
\usepackage{fontspec}
\usepackage{polyglossia}
\setdefaultlanguage[variant=mexican]{spanish}
% Los nombres chinos van como «pinyin (original)»: xeCJK da a los caracteres
% Han su propia fuente; sin ella XeLaTeX los omite con un aviso.
% FandolSong no cubre algunos caracteres de nombres (U+73FA); WenQuanYi Zen Hei sí.
\usepackage[AutoFallBack=true]{xeCJK}
\setCJKmainfont{FandolSong-Regular.otf}
\setCJKfallbackfamilyfont{\CJKrmdefault}{WenQuanYi Zen Hei}
% polyglossia no traduce los nombres de listings.
\AtBeginDocument{\providecommand{\lstlistingname}{}\renewcommand{\lstlistingname}{Listado}%
  \providecommand{\lstlistlistingname}{}\renewcommand{\lstlistlistingname}{Índice de listados}}
\usepackage{amsmath, amssymb}
\usepackage{graphicx}
\usepackage[margin=2.5cm]{geometry}
\usepackage[most]{tcolorbox}
\usepackage{etoolbox}
\usepackage{listings}
\usepackage{hyperref}
\usepackage{booktabs}
\usepackage{subcaption}
\usepackage{float}

% Highlight boxes
\newtcolorbox{knowledgebox}[1]{enhanced, colback=blue!5!white, colframe=blue!75!black, colbacktitle=blue!75!black, coltitle=white, fonttitle=\bfseries, title={#1}, attach boxed title to top left={yshift=-2mm, xshift=2mm}, boxrule=1pt, sharp corners}
\newtcolorbox{importantbox}[1]{enhanced, colback=yellow!10!white, colframe=yellow!80!black, colbacktitle=yellow!80!black, coltitle=black, fonttitle=\bfseries, title={#1}, sharp corners}
\newtcolorbox{warningbox}[1]{enhanced, colback=red!5!white, colframe=red!75!black, colbacktitle=red!75!black, coltitle=white, fonttitle=\bfseries, title={#1}, sharp corners}

% Listings
\lstset{language=Python, basicstyle=\ttfamily\small, keywordstyle=\color{blue}, stringstyle=\color{red!60!black}, commentstyle=\color{green!60!black}, breaklines=true, frame=single, numbers=left, numberstyle=\tiny\color{gray}, captionpos=b, extendedchars=true}

% Front-page metadata
\newcommand{\notetitle}{{[}LLM Agents F24{]} AI Agents for Enterprise Workflows — Nicolas Chapados \& Alexandre Lacoste}
\newcommand{\noteauthors}{Elaboradas a partir de materiales públicos del curso}
\newcommand{\notedate}{21 de octubre de 2024}
\newcommand{\videochannel}{Berkeley RDI}
\newcommand{\videopublishdate}{21 de octubre de 2024}
\newcommand{\videoduration}{aproximadamente 80 minutos}
\newcommand{\videourl}{https://www.youtube.com/live/-yf-e-9FvOc}
\newcommand{\videocoverpath}{cover.jpg}
\newcommand{\repourl}{https://github.com/hqhq1025/ai-course-notes}

% Footer with repo info
\usepackage{fancyhdr}
\pagestyle{fancy}
\fancyhf{}
\fancyfoot[C]{\thepage}
\fancyfoot[R]{\footnotesize\color{gray}\href{\repourl}{AI Course Notes}}
\renewcommand{\headrulewidth}{0pt}
\renewcommand{\footrulewidth}{0pt}

\begin{document}
\begin{titlepage}\centering
{\Large Notas de entrevista\par}\vspace{1.2cm}
{\huge\bfseries \notetitle\par}\vspace{0.8cm}
{\large \noteauthors\par}\vspace{0.3cm}
{\large \notedate\par}\vspace{1.2cm}
\ifdefempty{\videocoverpath}{{\small Por favor, incluya la ruta de la imagen de portada.\par}}{\includegraphics[width=0.82\textwidth,height=0.45\textheight,keepaspectratio]{\videocoverpath}\par}
\vfill
\begin{tcolorbox}[width=0.9\textwidth, colback=black!2!white, colframe=black!60, sharp corners]
\textbf{Autor o canal del video}: \videochannel\par
\textbf{Fecha de publicación}: \videopublishdate\par
\textbf{Duración del video}: \videoduration\par
\ifdefempty{\videourl}{}{\textbf{Enlace del video}: \href{\videourl}{\nolinkurl{\videourl}}\par}\par
\textbf{Página del proyecto}: \href{\repourl}{\nolinkurl{\repourl}}\par
\end{tcolorbox}
\end{titlepage}

\tableofcontents\newpage

\section{Introducción: la misión de los Agentes empresariales}
La sexta clase de Berkeley LLM Agents F24 estuvo a cargo de Nicolas Chapados y Alexandre Lacoste, de ServiceNow. Ambos ponentes subrayaron desde el inicio que ServiceNow es una empresa global con más de 25,000 empleados, cuyo producto central es el workflow. El valor de los Agentes no reside únicamente en la capacidad del modelo, sino en \textbf{«entregar resultados en el contexto real del trabajo del empleado»}, idea que se repite a lo largo del fragmento comprendido entre 00:01:11 y 00:01:37.

\begin{knowledgebox}{Los Agentes empresariales no pueden «operar fuera de contexto»}
La plataforma de ServiceNow conecta ITSM, recursos humanos, servicio al cliente y otras verticales; todo Agente debe primero tener claro «quién es esta persona», «qué está haciendo en este momento» y «qué ayuda espera del Agente» antes de diseñarse, a fin de cumplir con los requisitos de auditoría y seguridad.
\end{knowledgebox}

\subsection{Resumen de la sección}
En los Agentes de nivel empresarial, primero está la necesidad y después el modelo; dar prioridad al rol del usuario y al contexto del workflow es la primera lección para su implementación.

\section{De John a Sandy: cada paso del flujo de trabajo}
\subsection{Un ticket ordinario: la larga cadena del proceso manual}
En la historia de 00:02:11 a 00:11:21, John derrama café sobre su laptop por la mañana y abre un ticket con TI para solicitar un reemplazo del equipo; el ticket se asigna a Sandy, quien revisa uno por uno la base de conocimiento y los incidentes históricos, ajusta los permisos y deja un registro claro. Aunque ya está desplegada la IA generativa más reciente, Sandy sigue realizando una gran cantidad de pasos manuales, entre ellos verificar el control de acceso, invocar los sistemas relacionados y escribir el resumen de la solución.

\begin{warningbox}{La brecha del «dominó» de la automatización}
Aunque exista un resumen generado por IA, en los flujos de trabajo empresariales todavía hay cientos de escenarios que requieren criterio humano, verificación de permisos, evaluación de cumplimiento y coordinación entre varios sistemas; cualquier automatización debe conservar en este proceso la trazabilidad y un margen para la intervención humana.
\end{warningbox}

\subsection{Tareas «grano de arena»: poco frecuentes pero cruciales}
El presentador señala que la mayoría de las tareas de «bajo valor, baja frecuencia» son demasiado singulares para que los scripts tradicionales las cubran; estas «tareas grano de arena» se repiten millones de veces al día en las empresas de todo el mundo, pero cada una exige una respuesta rápida.

\begin{importantbox}{El alcance del Agente no debe medirse solo por la «alta frecuencia»
}
El Agente debe aprovechar las oportunidades de ser «explicable, trazable y ajustable con granularidad fina»: aunque una sola tarea tome poco tiempo, lo que realmente importa es el incremento de eficiencia y la satisfacción de los empleados que se generan a escala aggregate.
\end{importantbox}

\subsection{Resumen de la sección}
El ejemplo de John y Sandy muestra el valor potencial y los riesgos del Agente en escenarios con múltiples sistemas, roles y pasos, y sienta las bases para los requisitos de arquitectura y auditoría que se tratan más adelante.
\section{Los límites de rol entre el API Agent y el Web Agent}
\subsection{API Agent: basta una especificación clara para automatizar}
En el segmento de 00:08:35 a 00:09:35, el ponente describe al API Agent como «un conjunto de API o endpoints con nombre que intercambian principalmente información textual»; su conjunto de herramientas puede describirse con precisión mediante una especificación formal, lo cual resulta adecuado para flujos que ya cuentan con una interfaz de servicio.

\begin{knowledgebox}{El beneficio central del API Agent}
Cuando el conjunto de herramientas y los formatos de entrada y salida son claros y verificables, los API Agents pueden mantener una ejecución «de alta precisión y predecible», al mismo tiempo que conservan la reproducibilidad del prompt-engineering.
\end{knowledgebox}

\subsection{Web Agent: requiere percepción y razonamiento sobre la UI}
Los Web Agents se describen como agentes que «hacen clic, llenan formularios y navegan entre páginas dentro del navegador igual que una persona», pero al mismo tiempo cargan con la tarea de entender el objetivo, percibir la situación en el momento y planear a largo plazo (00:32:49--00:34:14).

\begin{warningbox}{El triple desafío del Web Agent}
Tiene que entender la intención del usuario, identificar la UI actual y además planear cómo llegar al resultado final; pero el punto final casi nunca está en la página actual, la observación presente por sí sola no permite determinar si la tarea ya se completó, y un error de ejecución puede provocar que se repitan operaciones o que el proceso se interrumpa.
\end{warningbox}

\subsection{Resumen de la sección}
Los API Agents y los Web Agents no se excluyen entre sí, sino que se complementan: los primeros garantizan estabilidad mediante la especificación, mientras que los segundos logran generalidad mediante la percepción; al diseñar un sistema hay que combinarlos de manera razonable según la tarea objetivo.

\section{TapeAgents: trazas de ejecución auditables}
Chapados presenta TapeAgents entre las 00:18:56 y las 00:20:52; se trata de un marco de trabajo recién liberado como código abierto cuyo objetivo es «tender un puente entre las herramientas de ingeniería y las herramientas de optimización». «Tape» es la abstracción unificada: en este registro se anotan el razonamiento, las acciones y las observaciones de todos los Agentes, por lo que resulta auditable y además puede ser leído por Agentes posteriores.

\begin{importantbox}{El «Best of Both Worlds» de TapeAgents}
«Queremos herramientas de desarrollo y depuración, pero también capacidad de optimización de prompts y de fine-tuning»; TapeAgents ofrece componentes de grano fino, flujos paralelos y runtime orchestration, y a la vez usa el tape como activo de datos para convertir el registro de ejecución en material de optimización.
\end{importantbox}

\subsection{Estructura y reutilización del tape}
El tape de TapeAgents no es un simple registro: contiene metadata estructurada y abundante que puede ser consumida por otros Agentes o por optimizadores (00:21:10--00:22:25). Varios Agentes pueden compartir el mismo tape; el Agente B lee el historial completo y elige el node que desea ejecutar, con lo que se construye un flujo de control modular.

\begin{knowledgebox}{El tape como «datos como servicio»}
El tape registra cada prompt, cada retroalimentación del entorno y cada instrucción de acción, lo que permite impulsar la depuración, la visualización, la optimización de prompts y la destilación maestro-alumno, formando así un ciclo cerrado que va del execution trace a los datos de fine-tune.
\end{knowledgebox}

\subsection{Resumen de la sección}
Gracias a un tape unificado con metadata abundante, TapeAgents logra la sinergia entre la colaboración entre varios Agentes, la depuración a nivel de ingeniería y la optimización automatizada, resolviendo la desconexión entre el «prompt manual» y la «optimización experimental».

\section{Optimización y GRATE: de modelos grandes a modelos compactos}
\subsection{Pipeline de entrenamiento con dos agentes y una métrica de calidad}
En el caso de 00:24:31--00:31:15, ServiceNow primero usa el modelo grande Llama 3.5, con 405B de parámetros, para producir interacciones de alta calidad, y después usa los registros de tape como datos para hacer fine-tuning de un modelo Llama de apenas 8B de parámetros y así recuperar el desempeño. El resultado es que el modelo ajustado obtiene una puntuación GRATE equivalente a la del modelo grande original, pero logra un ahorro de más de 300x en el costo de operación.

\begin{knowledgebox}{Puntuación de calidad GRATE}
GRATE es un conjunto de atributos de calidad cuantificables (grounded, responsive, accurate, transparent, effective); cada dimensión puede medirse y optimizarse de manera independiente, lo que facilita usar el tape como señal de reward o de regresión durante el entrenamiento.
\end{knowledgebox}

\subsection{Compromiso entre costo y desempeño}
La gráfica de métricas de 00:30:25--00:31:38 muestra: el eje x representa el costo por millón de turnos de conversación, y el eje y representa la puntuación GRATE; los modelos grandes siguen en la esquina superior derecha, pero los modelos ajustados generados mediante tape pueden alcanzar una calidad cercana con apenas una milésima parte del costo.

\subsection{Resumen de la sección}
TapeAgents no solo facilita el desarrollo con ingeniería, sino que también nos permite convertir un agente «maestro» costoso en un «agente estudiante» de alta calidad y bajo costo; la curva de costo y desempeño puede visualizarse mediante GRATE.

\section{Práctica y retos del Web Agent}
\subsection{Planear un viaje a la GTC con un Web Agent}
En la demostración de 00:15:41 a 00:16:39, el Agent primero le pregunta al usuario «quiero ir a GTC 2024», después visita sucesivamente Google y Google Maps, y finalmente encuentra San Jose Convention Center; luego deja que el navegador ejecute el formulario y la navegación de la ruta. Este fragmento muestra que el Web Agent puede completar la tarea apoyándose en la interacción con la página incluso sin disponer de una API, y también muestra que necesita preguntarle constantemente al ser humano y hacer replanning dinámico.
\footnote{El intervalo de video correspondiente a este fragmento es 00:15:41--00:16:39.}

\begin{importantbox}{Los límites de la colaboración entre el humano y el Web Agent}
En los escenarios donde no se dispone de una API, el Web Agent depende sobre todo de la percepción visual y de un ciclo de instrucciones humanas: tiene que extraer el DOM o el texto del navegador, luego razonar, luego ejecutar clics o entradas de texto, y durante ese proceso pedirle información adicional o confirmación al ser humano; esto es la manifestación de que «el modo es la capacidad».
\end{importantbox}

\subsection{Planificación a largo plazo y limitaciones de observación}
El punto final del Web Agent casi nunca está en la página actual, por lo que debe mantener una estrategia interna para evitar descubrir, tras varios clics, que tomó un camino equivocado. Además, el SRT de 00:33:54 a 00:34:14 señala que, incluso cuando el Agent elabora un plan, también debe garantizar la precisión de las acciones que ejecuta, o de lo contrario tenderá a repetir pasos.

\begin{warningbox}{El riesgo de separar Planning y Execution}
El Web Agent puede elaborar un plan excelente, pero fallar porque la interfaz cambió, el texto del botón es ligeramente distinto o se requiere una verificación de varios niveles; en ese caso hay que retroalimentar el error hacia el tape, para que los nodos posteriores aprendan una manera de ejecución más robusta.
\end{warningbox}

\subsection{Resumen de la sección}
El Web Agent es el resultado combinado de que «el humano esté en el loop» y de que «el Agent esté en el loop»; la semántica de la interfaz de usuario, la planificación a largo plazo y la precisión en la ejecución de las acciones necesitan mecanismos sistemáticos que las garanticen.

\section{Consideraciones de gobernanza y despliegue}
En el segmento de 00:03:48--00:04:02, el ponente enfatizó la importancia de la confianza, la gobernanza y la seguridad, que son la condición previa para que un Agente de IA pueda ser autorizado en un entorno empresarial. El tape mismo de TapeAgents también constituye una cadena de datos auditable, que sirve de base para el cumplimiento normativo, la depuración y la reversión.

\begin{warningbox}{El cumplimiento no es «documentación», sino «proceso»}
Desplegar un Agent requiere definir con claridad «quién aprueba», «quién puede modificar el prompt» y «quién recibe el resumen de ejecución»; de lo contrario, aunque el modelo tenga un buen desempeño, el equipo de seguridad puede darlo de baja por la falta de una cadena de gobernanza.
\end{warningbox}

\subsection{Resumen de la sección}
El Agent debe alinearse con el sistema de gobernanza de la empresa; la cadena de datos de TapeAgents ayuda a hacer transparente el proceso de ejecución y a facilitar la auditoría y la rendición de cuentas.
\section{Síntesis y ampliación}
\begin{tabular}{@{}p{4.5cm}p{11cm}@{}}
\toprule
Dimensión & Puntos clave \\
\midrule
Panorama del escenario & ServiceNow, mediante la historia de John/Sandy, revela oportunidades de automatización de tickets poco frecuentes pero críticas. \\
Ruta de arquitectura & El API Agent usa un conjunto de herramientas normalizado, el Web Agent opera el navegador; ambos deben complementarse. \\
TapeAgents & El tape es un registro estructurado que permite depurar, optimizar y destilar, y resuelve la separación entre el desarrollo del prompt y el del modelo. \\
Ciclo de optimización & Se usa un modelo grande para hacer demostraciones y el tape para producir datos con los que se hace fine-tuning de un modelo pequeño, lo que hace que la calidad de GRATE sea visible y el costo, controlable. \\
Requisitos de gobernanza & Auditoría de extremo a extremo, aprobación de cumplimiento normativo; el prompt y el resumen de ejecución deben poder consultarse. \\
\bottomrule
\end{tabular}

\subsection{Lecturas adicionales}
\begin{itemize}
    \item El informe técnico oficial de TapeAgents y los recursos de GitHub (el presentador da explícitamente el código QR y el enlace en 00:19:05).\newline
    \item El artículo original de WorkArena/BrowserGym: ofrece una evaluación estandarizada y un benchmark para Web Agents de nivel empresarial.
    \item La presentación de la «Enterprise Workflow Platform» de ServiceNow ayuda a entender las prácticas estándar de colaboración entre múltiples sistemas y de gobernanza.
\end{itemize}

\end{document}
